\section{Introduction}
\label{sec:intro}

Visual attention is the set of mechanisms by which a behaving animal commits its
limited processing resources to the small fraction of the visual scene that is
momentarily most relevant to its goals \citep{carrasco2011visual,
desimone1995neural, reynolds2009normalization, rust2022priority}. Decades of
psychophysics and neurophysiology have established a robust catalogue of its
behavioural and neural signatures: faster and more accurate responses to attended
stimuli \citep{posner1980attention, carrasco2011visual, hawkins1990visual}, the
graded dependence of those benefits on the reliability and value of predictive
cues \citep{mayo2016graded, eckstein2013visual}, and the modulation of single-unit
responses and inter-areal synchrony throughout the visual hierarchy
\citep{reynolds2000attention, treue1999reshaping, fries2001modulation,
gregoriou2009high}. Yet a unified computational account of \emph{why} attention is
organised the way it is --- of which structural features of the underlying
circuitry are necessary for attentive behaviour to arise --- remains incomplete
\citep{mehrani2023self, hassanin2024visual, lindsay2020attention}.

\paragraph{Attention as two coupled computations.}
A useful way to frame the problem is to recognise that an attending animal must
perform two logically distinct computations concurrently. It must form a
\emph{policy} --- a moment-by-moment decision about what to do, such as whether to
report that a behaviourally relevant event has occurred --- and it must form a
\emph{value estimate} --- a prediction of how good its current situation is, which
under any reward-driven account is precisely the quantity that supplies the
teaching signal for the policy \citep{sutton2018reinforcement, dayan2008decision,
niv2009reinforcement}. The two computations draw on overlapping but distinguishable
information. \emph{When and where to act} is driven by bottom-up, image-grounded
evidence accumulated toward a commitment threshold
\citep{gold2007neural, roitman2002response, ding2011neural}; \emph{how valuable the
moment is} depends on accumulated context --- which location was cued, how reliable
the cue is, how much reward is at stake, how much time has elapsed
\citep{platt1999neural, sugrue2004matching, stanisor2013neural}. This division maps
onto a deep distinction in the biology, between a sensory-grounded priority and
commitment system and a value/prediction-error system, that we take as the central
motivation for the present work.

\paragraph{The priority/commitment system.}
A network of topographically organised areas --- the lateral intraparietal area
(LIP), the frontal eye field (FEF), and the superior colliculus (SC) --- maintains
a ``priority map'' that combines bottom-up conspicuity with top-down task relevance
and is read out to guide both overt orienting and covert attention
\citep{bisley2010lip, fecteau2006salience, thompson2005neuronal,
veale2017how}. The activity of these areas carries the signature of an evidence
integrator: LIP firing ramps as the running time-integral of momentary sensory
evidence toward a decision bound \citep{roitman2002response, huk2005neural}, and
FEF and SC convert that accumulation into a committed response
\citep{ding2011neural, stine2023differential}. Causal manipulations place these
areas upstream of attentive behaviour rather than merely correlated with it:
sub-threshold microstimulation of FEF or SC \emph{enhances} covert perceptual
sensitivity at the corresponding location without an eye movement
\citep{moore2001control, moore2003selective, cavanaugh2006subcortical,
muller2005microstimulation}, while their inactivation \emph{impairs} covert
selection \citep{wardak2006contribution, lovejoy2010inactivation,
zenon2012attention, bollimunta2018fefsc}. Critically, the contribution of the
image-grounded signal is dissociable from the value it carries: in FEF the
attentional effect is borne by visually responsive cells, with pure
movement-related cells suppressed during covert attention
\citep{thompson2005neuronal, gregoriou2012lesions}, and in SC the deficit produced
by inactivation survives intact cortical sensory gain \citep{zenon2012attention}.

\paragraph{The value/actor--critic system.}
Running in parallel, the cortico-basal-ganglia loop implements a reinforcement-
learning actor--critic: action selection is associated with dorsal striatum (the
``actor''), while value prediction and the reward-prediction-error teaching signal
are associated with ventral striatum and midbrain dopamine (the ``critic'')
\citep{houk1995model, joel2002actor, odoherty2004dissociable, maia2009reinforcement,
collins2014opponent}. The cornerstone of this account is the close correspondence
between phasic dopamine firing and the temporal-difference reward-prediction error
\citep{schultz1997neural, hollerman1998dopamine, montague1996framework}, together
with the dopamine dependence of cortico-striatal plasticity
\citep{reynolds2002dopamine, morita2014striatal}. The two systems are not
independent: the actor learns \emph{only} through a three-factor plasticity rule in
which the critic's dopaminergic error is the gating third factor
\citep{kusmierz2017learning, gerstner2018eligibility}, so the value system and the
action system are bound into a single closed learning loop. This furnishes a strong
prior expectation --- that policy and value computations cannot be severed without
breaking learning --- that our model makes precise and tests.

\paragraph{Existing computational models leave the organisational question open.}
Models of visual attention have captured many of its features one or a few at a
time --- normalisation-based gain control \citep{reynolds2009normalization},
bottom-up saliency maps \citep{itti2001computational, koch1985shifts}, hierarchical
selection \citep{tsotsos1995modeling}, guided visual search
\citep{wolfe2021guided}, and predictive/active-inference accounts
\citep{feldman2010attention, parr2017working, mirza2018human}. Few learn their
attentional strategy from reward as animals do
\citep{chelazzi2013reward, anderson2016attention, failing2018selection,
mnih2014recurrent}, and fewer still expose an interpretable, causally manipulable
internal attention signal. Contemporary transformers, whose self-attention
mechanism superficially resembles biological selection
\citep{vaswani2017attention, dosovitskiy2020image, khan2022transformers}, in
practice emphasise low-level grouping rather than task-dependent, goal-driven
selection \citep{mehrani2023self}, and lack the recurrent internal state and
capacity limits that are central to biological attention and working memory
\citep{luck1997capacity, awh2006interactions, panichello2021attvwm}. Recurrent
neural networks supply the requisite dynamics \citep{hochreiter1997long,
beck2024xlstm, sussillo2014neural, mante2013context} but are not, by default,
interpretable as attention maps. None of these traditions has asked the specific
question we pose here: holding the learning signal and the components fixed, does
the \emph{wiring} between the policy and value computations determine whether
attentive behaviour is acquired?

\paragraph{This work.}
We developed a recurrent vision transformer that integrates spatial working memory
with self-attention into a ``recurrent self-attention'' mechanism and is trained
end-to-end by reinforcement learning on sparse, animal-like reward
\citep{springenberg2024offline, bellemare2017distributional}. Building on a
single-stream version of this model that already reproduces hallmark behavioural
and causal signatures of primate covert attention, we partition the agent into two
specialised recurrent cross-attention pathways: a \emph{$\mu$ pathway} that queries
the live image and drives the actor (the decision of when to declare a change), and
a \emph{$q$ pathway} that reads accumulated memory and drives the distributional
critic (the value estimate). We then manipulate a single structural variable ---
inter-pathway coupling --- in the spirit of a lesion experiment, comparing (i) a
\emph{shared-memory} design in which the pathways cross-talk through a common
recurrent state, (ii) an \emph{independent} design in which each pathway is given a
private memory it never shares, and (iii) a \emph{swapped-routing} design that keeps
the coupling but lets the memory pathway command the policy. We report four central
findings. First, coupling is required: only the shared-memory design learns the task,
and the independent design fails to acquire it at all. Second, the requirement is
asymmetric: the image-grounded pathway must drive the policy; routing the decision
through the memory pathway collapses the agent identically. Third, the coupled
architecture is general --- it scales to a nine-stimulus set-size manipulation and to
an environment with competing distractors, reproducing the attentional signatures in
each. Fourth, the trained model yields concrete, falsifiable predictions for systems
neuroscience: a representational dissociation (change-evidence should be decodable
from the policy pathway, value/context from the value pathway) and a causal
dissociation (perturbing the value pathway should spare detection while degrading
value-based modulation, whereas perturbing the policy pathway should abolish
detection). We use the deliberately interpretable, manipulable structure of the
model not as an end in itself but as a tool: a generative source of mechanistic
hypotheses that the experimental study of LIP, FEF, SC and the basal ganglia can
test.
